% Options for packages loaded elsewhere
\PassOptionsToPackage{unicode}{hyperref}
\PassOptionsToPackage{hyphens}{url}
\documentclass[
  12pt,
]{article}
\usepackage{xcolor}
\usepackage{amsmath,amssymb}
\setcounter{secnumdepth}{-\maxdimen} % remove section numbering
\usepackage{iftex}
\ifPDFTeX
  \usepackage[T1]{fontenc}
  \usepackage[utf8]{inputenc}
  \usepackage{textcomp} % provide euro and other symbols
\else % if luatex or xetex
  \usepackage{unicode-math} % this also loads fontspec
  \defaultfontfeatures{Scale=MatchLowercase}
  \defaultfontfeatures[\rmfamily]{Ligatures=TeX,Scale=1}
\fi
\usepackage{lmodern}
\ifPDFTeX\else
  % xetex/luatex font selection
\fi
% Use upquote if available, for straight quotes in verbatim environments
\IfFileExists{upquote.sty}{\usepackage{upquote}}{}
\IfFileExists{microtype.sty}{% use microtype if available
  \usepackage[]{microtype}
  \UseMicrotypeSet[protrusion]{basicmath} % disable protrusion for tt fonts
}{}
\makeatletter
\@ifundefined{KOMAClassName}{% if non-KOMA class
  \IfFileExists{parskip.sty}{%
    \usepackage{parskip}
  }{% else
    \setlength{\parindent}{0pt}
    \setlength{\parskip}{6pt plus 2pt minus 1pt}}
}{% if KOMA class
  \KOMAoptions{parskip=half}}
\makeatother
\usepackage{longtable,booktabs,array}
\usepackage{calc} % for calculating minipage widths
% Correct order of tables after \paragraph or \subparagraph
\usepackage{etoolbox}
\makeatletter
\patchcmd\longtable{\par}{\if@noskipsec\mbox{}\fi\par}{}{}
\makeatother
% Allow footnotes in longtable head/foot
\IfFileExists{footnotehyper.sty}{\usepackage{footnotehyper}}{\usepackage{footnote}}
\makesavenoteenv{longtable}
\setlength{\emergencystretch}{3em} % prevent overfull lines
\providecommand{\tightlist}{%
  \setlength{\itemsep}{0pt}\setlength{\parskip}{0pt}}
\usepackage{bookmark}
\IfFileExists{xurl.sty}{\usepackage{xurl}}{} % add URL line breaks if available
\urlstyle{same}
\hypersetup{
  hidelinks,
  pdfcreator={LaTeX via pandoc}}

\author{SCX}
\date{2026}

\begin{document}

\section{Proposition 6: Error-Driven State
Discovery}\label{proposition-6-error-driven-state-discovery}

\begin{quote}
The optimal state partition is not pre-specified by humans, but is automatically discovered in the error-relevant feature subspace. States should be defined by ``where the model fails,'' not by ``what humans think is important.''
\end{quote}

\begin{center}\rule{0.5\linewidth}{0.5pt}\end{center}

\subsection{1 Statement}\label{statement}

\subsubsection{1.1 Formal Statement}\label{formal-statement-prop6}

The optimal state partition \(\mathcal{S}^*\)
is not pre-specified by humans, but is automatically discovered in the error-relevant feature subspace.

Specifically, given a Layer 1 encoding \(\varphi_1\) and residual vector
\(r\), states should be clustered in the subset of \(\varphi_1\) dimensions most correlated with \(r\):

\[\mathcal{S}^* = \text{Cluster}(\varphi_1[:, \text{top\_dims}], K)\]

where:

\[\text{top\_dims} = \text{argsort}(\text{MI}(\varphi_1[:, j], r))\]

i.e., select the subset of \(\varphi_1\) dimensions with the highest mutual information with residual \(r\) as the clustering space.

\subsubsection{1.2 Two-Layer Descriptor System}\label{two-layer-descriptor-system}

\textbf{Layer 1 descriptors (provided by researchers)}: The material's ``identity label'' ---
element type, stoichiometric ratio, crystal structure ---
temperature/pressure range, structure type (bulk/surface/defect/interface) ---
these descriptors define the \textbf{initial sampling space}

\textbf{Layer 2 descriptors (discovered by algorithm)}: Failure modes revealed by error analysis ---
coordination number (CN), maximum bond stretch, local strain ---
sp2/sp3 ratio, atomic displacement norm, local coordination changes ---
ACE basis residual components, angular distribution of force errors ---
these descriptors are \textbf{not manually specified by researchers}---they are automatically computed, filtered, and combined by the algorithm during error analysis

\subsubsection{1.3 Core Idea}\label{core-idea-prop6}

\begin{quote}
Human-defined descriptors may fail to capture the model's true failure modes.

States should be defined by ``where the model fails,'' not by ``what humans think is important.''
\end{quote}

\begin{center}\rule{0.5\linewidth}{0.5pt}\end{center}

\subsection{2 Formal Statement}\label{formal-statement-detail}

\subsubsection{2.1 Mathematical Formalization}\label{mathematical-formalization}

Let \(\varphi_1(x)\) be the Layer 1 encoding (e.g., the embedding output of SCXStateEncoder),
\(r(x) = |f_{\text{model}}(x) - f_{\text{true}}(x)|\)
be the residual.

Define the error-relevant subspace:

\[D^* = \{j: \text{MI}(\varphi_1[:, j], r) > \tau\}\]

i.e., among the dimensions of the Layer 1 encoding, the set of dimension indices whose mutual information with residual \(r\) exceeds threshold \(\tau\).

Then the clustering \(\mathcal{S}^*\) on \(D^*\) satisfies:

\[\mathbb{E}[\text{Var}(r \mid s \in \mathcal{S}^*)] \leq \mathbb{E}[\text{Var}(r \mid s \in \mathcal{S}_{\text{human}})]\]

i.e., the within-group residual variance of error-driven states does not exceed the within-group residual variance of human-defined states.

\subsubsection{2.2 Algorithm Workflow}\label{algorithm-workflow}

\textbf{Phase A: Initial Training}

\begin{verbatim}
Input: Initial DFT dataset D_0 defined by Layer 1 descriptors
Process: Train Single ACE baseline M_0
Output: M_0, test set error distribution
\end{verbatim}

\textbf{Phase B: Error Landscape Analysis}

\begin{verbatim}
Input: M_0, test set
Process:
  1. Compute per-atom force error and per-structure error for each test structure
  2. Compute candidate Layer 2 descriptors for each atom:
     - CN (at different cutoff radii), max/average bond stretch
     - Local strain tensor invariants, bond angle distribution statistics
     - Voronoi volume, ACE basis component contributions
  3. Filter descriptors most correlated with error using MI / SHAP / permutation importance
  4. UMAP/PCA dimensionality reduction in the selected descriptor space
  5. HDBSCAN clustering to discover error clusters
Output:
  - Error cluster map (cluster_id → descriptor range → average error → sample count)
  - Definition of high-error regions (convex hull or interval in descriptor space)
\end{verbatim}

\textbf{Phase C: Structure Proposal}

\begin{verbatim}
Input: High-error region definitions
Process:
  1. Generate candidate structures within high-error descriptor regions
  2. Deduplicate (compare structural fingerprints with existing training set)
  3. Rank (prioritize structures with largest error, highest information content)
Output: Structure proposal file + expected error improvement estimate
\end{verbatim}

\textbf{Phase D: Iterative Enhancement}

\begin{verbatim}
Input: New DFT data D_new
Process:
  1. Merge D = D_old ∪ D_new
  2. Retrain or incrementally train
  3. Return to Phase B to reassess error landscape
  4. If new high-error clusters appear → Phase C
  5. If error landscape flattens → terminate
\end{verbatim}

\textbf{Phase E (Optional): Expertization}

\begin{verbatim}
Input: Final error cluster definitions
Process:
  1. Each error cluster automatically defines an expert domain
  2. Train residual expert for each domain
  3. Manage expert modules with Expert Algebra
  4. Export as deployable potential function
\end{verbatim}

\begin{center}\rule{0.5\linewidth}{0.5pt}\end{center}

\subsection{3 Theorem: Error-Driven States Dominate Human-Defined
States}\label{theorem-error-driven-states-dominate-human-defined-states}

\textbf{Theorem 6.1 (Advantage of Error-Driven States)}: Let
\(\mathcal{S}_{\text{human}}\)
be a state partition defined by human prior knowledge (e.g., classification by crystal structure, chemical composition, etc.), and \(\mathcal{S}^*\)
be the state partition discovered via error-driven discovery. Under the same number of states \(K = |\mathcal{S}|\):

\[\mathbb{E}_{s \sim \mathcal{S}^*}[\text{Var}(r \mid s)] \leq \mathbb{E}_{s \sim \mathcal{S}_{\text{human}}}[\text{Var}(r \mid s)]\]

\textbf{Proof sketch}: By the data processing inequality of mutual information, the error-relevant subspace \(D^*\)
retains the most information about \(r\). The clustering objective of K-means or HDBSCAN on \(D^*\)
is equivalent to minimizing within-group variance:

\[\min_{\mathcal{S}} \sum_{s \in \mathcal{S}} \sum_{x \in s} \|\varphi_1(x)[D^*] - \mu_s\|^2\]

Whereas \(\mathcal{S}_{\text{human}}\) is defined on the original space \(\varphi_1\) (which may contain dimensions unrelated to
\(r\)). Therefore \(\mathcal{S}^*\) has conditional variance on \(r\) that is at least no worse than \(\mathcal{S}_{\text{human}}\). \(\square\)

\begin{center}\rule{0.5\linewidth}{0.5pt}\end{center}

\subsection{4 Connection to SCX Core}\label{connection-to-scx-core}

\subsubsection{4.1 Architecture Mapping}\label{architecture-mapping}

\begin{itemize}
\tightlist
\item
  \textbf{Layer 1} = SCXStateEncoder.encode() --- initial embedding mapping
\item
  \textbf{Layer 2} = ErrorDrivenEncoder.fit\_error\_states() ---
  error-conditioned state discovery
\item
  \textbf{Two-layer integration} = TwoLayerStateDiscovery --- the complete automated state discovery pipeline
\end{itemize}

\subsubsection{4.2 Why a Single-Layer Encoder is Insufficient}\label{why-single-layer-insufficient}

Why a single-layer encoder (e.g., 12-dimensional MLIP encoder) fails:

\begin{enumerate}
\def\labelenumi{\arabic{enumi}.}
\tightlist
\item
  \textbf{Generic embeddings are error-agnostic}: Pre-trained embeddings may encode extensive global structural information unrelated to expert failure modes
\item
  \textbf{Fixed dimensionality has limited capacity}: 12-dimensional embeddings may not simultaneously encode chemical environment information and error pattern information
\item
  \textbf{Lack of error feedback}: The error signal is not backpropagated into the state construction process
\end{enumerate}

The two-layer descriptor system resolves these issues through error-conditioned projection.

\subsubsection{4.3 Relationship to the Information Bottleneck}\label{relationship-info-bottleneck}

Two-layer state discovery corresponds to an iterative solution of the \textbf{conditional information bottleneck}:

\[\min_{\Pi} I(\varphi_1; \tilde{X} \mid \mathcal{F}) - \beta I(\tilde{X}; r \mid \mathcal{F})\]

where \(\mathcal{F} = \{f_1, \dots, f_M\}\) is the expert set, \(r\)
is the residual, and \(\tilde{X}\) is the state representation.

\begin{center}\rule{0.5\linewidth}{0.5pt}\end{center}

\subsection{5 Implications}\label{implications}

\begin{enumerate}
\def\labelenumi{\arabic{enumi}.}
\item
  \textbf{State partitions are not fixed}---they evolve as the model improves. As the model error landscape flattens, state boundaries need to be redefined.
\item
  \textbf{Layer 2 descriptors = error-conditioned subset of Layer 1 features}---automatically filter out the feature dimensions most relevant to model failure, excluding irrelevant ``background'' dimensions.
\item
  \textbf{This distinguishes SCX from all human-feature-based clustering methods}:

  \begin{itemize}
  \tightlist
  \item
    Traditional approach: researchers define states (``bulk'', ``surface'',
    ``defect''), but these may not reflect the model's true failure modes
  \item
    SCX + two-layer descriptors: the algorithm automatically discovers states from the error distribution
  \end{itemize}
\item
  \textbf{Emergence of expert partitioning}: Expert partitions should not be empirically defined by humans, but should emerge from the error distribution via the algorithm.
\item
  \textbf{Integration with Expert Governance Protocol}:

\begin{verbatim}
Two-layer descriptor framework (upper layer: auto-discover expert boundaries)
     │
     ▼
Expert Algebra framework (lower layer: manage expert modules)
\end{verbatim}
\end{enumerate}

\begin{center}\rule{0.5\linewidth}{0.5pt}\end{center}

\subsection{6 Key Validation Experiments}\label{key-validation-experiments}

\begin{longtable}[]{@{}p{0.25\linewidth}p{0.25\linewidth}p{0.42\linewidth}@{}}
\toprule\noalign{}
\begin{minipage}[b]{\linewidth}\raggedright
Experiment
\end{minipage} & \begin{minipage}[b]{\linewidth}\raggedright
Hypothesis
\end{minipage} & \begin{minipage}[b]{\linewidth}\raggedright
Success Criterion
\end{minipage} \\
\midrule\noalign{}
\endhead
\bottomrule\noalign{}
\endlastfoot
Error vs CN & CN<3.5 region error > CN≈4 region &
p<0.01, effect size>2x \\
Error vs bond stretch & stretch>10\% region error significantly larger &
p<0.01 \\
Descriptor separability & Different batches separable in descriptor space & silhouette score > 0.3 \\
Targeted vs uniform & error-guided sampling achieves same accuracy with fewer DFT & RMSE reduction > 20\% at equal budget \\
Auto expert vs manual expert & Algorithm-discovered expert partition no worse than manual & metrics comparable at equal expert count \\
\end{longtable}

\begin{center}\rule{0.5\linewidth}{0.5pt}\end{center}

\subsection{References}\label{references-prop6}

\begin{enumerate}
\def\labelenumi{\arabic{enumi}.}
\tightlist
\item
  EGP Two-Layer Descriptor Framework.
  \texttt{EGP\_TwoLayer\_Descriptors\_AlgorithmDriven\_Expert\_Partition\_2026-06-24.md}
\item
  Tishby, N., Pereira, F. C., \& Bialek, W. (1999). The information
  bottleneck method. \emph{Allerton Conference}.
\item
  Gauge-Normalized Residual ACE Expert Algebra Framework.
  \texttt{CodexKnowledge/}
\item
  SCX State-Conditioned eXpertise Core Framework.
  \texttt{01\_SCX\_Core\_Framework\_Math\_Analysis.md}
\end{enumerate}

\end{document}
